% Sec. V (continued): V-C Protocol Properties, V-D Timing and Liveness,
% V-E Griefing and the Wormhole.
\subsection{Protocol Properties}\label{sec:properties}

We state the guarantees of \code{lightning.spthy} as lemmas. Each is named after the \tamarin{}
property that establishes it. Adversarial capability is explicitly modelled as long-term key
compromise (\code{Compromise_Ltk}). In addition, two compromised parties on the same payment
can relay a preimage to each other around the parties between them (rule \code{AttackerSkip},
Sec.~\ref{sec:core-skip}). We apply honesty assumptions explicitly, in each lemma, not
globally; the exceptions are the restrictions of Sec.~\ref{sec:restrictions}, which hold in
every trace. Unless stated otherwise, every lemma of this section is proved in the default mode
of \code{lightning.spthy} (\texttt{make verify}), under all nine restrictions, including T1 to T3.

Table~\ref{tab:groups} sorts all lemmas into groups, one per stage of a payment that they
constrain. Each circled number in Figs.~\ref{fig:inflight} and~\ref{fig:wormhole} matches the
group that constrains that stage. We present the central lemmas below;
Appendix~\ref{app:catalogue} gives every statement.

\begin{table*}[t]
  \caption{Lemma groups of \texttt{lightning.spthy}; the circled numbers in
    Figs.~\ref{fig:inflight} and~\ref{fig:wormhole} refer to groups 1 to 7. A: all-traces;
    E: exists-trace. Unmarked rows: the 55 lemmas of \texttt{make verify}, all verified
    (32 all-traces safety lemmas, the 13 hop-bypass lemmas of group 9, 9 exists-trace lemmas
    found automatically, and \texttt{Wormhole\_Fee\_Theft\_Reachable}, replayed from a saved
    116-step proof). $\dagger$:~proved in \texttt{make general} (liveness restrictions removed;
    Sec.~\ref{sec:bypass}). $\ddagger$:~exists-trace, checked with T3 removed
    (Table~\ref{tab:necessity}).}
  \label{tab:groups}
  \centering
  \footnotesize
  \setlength{\tabcolsep}{4pt}
  \renewcommand{\arraystretch}{1.05}
  \begin{tabular}{@{}rL{2.75cm}L{13.6cm}c@{}}
    \toprule
    \# & Group & Lemmas & Kind\\
    \midrule
    0 & Channel layer, keys
      & \code{state_update}, \code{delayed_funds}, \code{instant_funds},
        \code{settlement_is_traceable}, \code{No_Punishment_Without_Cheating},
        \code{Funds_Locked_Before_Update}, \code{Update_Requires_Negotiation},
        \code{Ltk_Known_Implies_Compromised} & A\\
      & \quad non-vacuity & \code{Protocol_execution}, \code{Cooperative_Close_Execution} & E\\
    1 & Invoice \& preimage
      & \code{Invoice_Authenticates_Settlement}, \code{Preimage_Secret_Until_Released},
        \code{Preimage_Never_Known}, \code{Invoice_Has_Secret_Preimage},
        \code{Invoice_Released_Once} & A\\
      & \quad reachable & \code{Forged_Invoice_Requires_Key_Compromise} & E\\
    2 & Lock causality
      & \code{Forward_Requires_Incoming}, \code{HTLC_On_Opened_Channel},
        \code{Offer_Requires_Invoice} & A\\
    3 & Release causality
      & \code{Claim_Requires_Release}, \code{Settle_Requires_Receiver_Release},
        \code{Redeem_Requires_Receiver_Release}, \code{Fulfilled_Requires_Release},
        \code{Fulfilled_Source} & A\\
    4 & Atomicity
      & \code{Payment_Atomicity_Under_Liveness}, \code{Settle_Excludes_Sender_Refund},
        \code{OneOutcomePerHTLC} & A\\
      & \quad non-vacuity
      & \code{Multihop_Routing_Reachable}, \code{Sanity_Honest_Payment_Completes},
        \code{Sanity_Atomic_Payment_Completes}, \code{Distinct_Parties_Configuration} & E\\
    5 & Fees
      & \code{Fee_Only_On_Successful_Forward} (per hop, incoming amount = outgoing amount + fee
        by construction) & A\\
    6 & Refund / griefing & \code{Refund_Requires_Timeout}, \code{Loss_Requires_Inaction} & A\\
      & \quad reachable & \code{Refund_Possible} & E\\
    7 & Wormhole & \code{Skipped_Requires_Compromise} & A\\
      & \quad reachable
      & \code{Wormhole_Fee_Theft_Reachable} (saved 116-step proof, replayed) & E\\
    8 & Intermediary safety, timing
      & \code{Intermediary_Never_Loses_Under_Liveness}, \code{Intermediary_Never_Loses_Strong},
        \code{Timeout_Race_Blocked}, \code{ForwardTimeGap} & A\\
      & \quad non-vacuity & \code{Sanity_Compromised_Intermediary_Can_Lose} & E\\
      & \quad necessity $\ddagger$ & \code{Necessity_IMC_Honest_Intermediary_Loses} & E\\
    9 & Hop bypass
      & \code{First_Hop_Skip_Requires_Compromise}, \code{Second_Hop_Skip_Requires_Compromise},
        \code{Hop3_Skip_Requires_Compromise}, \code{Hop4_Skip_Requires_Compromise},
        \code{Hop5_Skip_Requires_Compromise}, \code{Hop6_Skip_Requires_Compromise} & A\\
      & \quad helpers
      & \code{Forward_Unique_Per_Incoming}, \code{SettledHop_From_Forward},
        \code{SettledHop_Uses_Outgoing}, \code{No_Settle_And_Refund_Same_Hop},
        \code{No_Fulfill_On_Forwarded_Incoming}, \code{Redeemed_Source},
        \code{SenderSettled_Uses_Fulfilled} & A\\
      & \quad any length $\dagger$ & \code{Any_Hop_Skip_Requires_Compromise} & A\\
      & \quad helpers $\dagger$
      & \code{Forward_Parent}, \code{Hop_Redeem_Propagates}, \code{Sender_Channel_Redeemed},
        \code{Root_First}, \code{Next_Is_Child}, \code{Route_Redeemed_Gen},
        \code{Route_Redeemed_No_Compromise} & A\\
    \bottomrule
  \end{tabular}
\end{table*}

The first guarantee is atomicity: on a route with two forwarders and no key compromise, once
the receiver has redeemed its HTLC, the sender can no longer be refunded.

\begin{lemma}[Payment Atomicity]\label{lem:atomicity}
For any two consecutive forwarders $F_1$ and $F_2$ on a payment with hash $y$ (the incoming HTLC
of $F_2$ is the outgoing HTLC of $F_1$): if the outgoing HTLC of $F_2$ on $y$ is redeemed and no
party's key was ever compromised, then the incoming HTLC of $F_1$ on $y$ is never refunded, by
anyone.
\end{lemma}
\begin{spthy}
lemma Payment_Atomicity_Under_Liveness:
  "All F1 F2 Prev1 Next1 Prev2 Next2 ptrSF1 ptrF1F2 ptrF2R y #i #j #k.
      HTLCForwarded(F1, Prev1, Next1, ptrSF1, ptrF1F2, y) @ #i
    & HTLCForwarded(F2, Prev2, Next2, ptrF1F2, ptrF2R, y) @ #j
    & Redeemed(ptrF2R, y) @ #k
    & (not Ex P #l. LtkCompromised(P) @ #l)
   ==> not (Ex Up #m. Refunded(Up, ptrSF1, y) @ #m)"
\end{spthy}
Since \code{Forward_HTLC} is one generic rule, $F_1$ and $F_2$ may sit anywhere on a route of
any length; on a three-hop route, \code{ptrSF1} is the sender's HTLC and \code{ptrF2R} the
receiver's. \code{Sanity_Atomic_Payment_Completes} shows that the premises are satisfiable.

The second guarantee protects the intermediaries, who lock their own funds for a payment that
is not theirs.

\begin{lemma}[Intermediary Safety]\label{lem:intermediary}
If a forwarder's outgoing HTLC is redeemed while its incoming HTLC is refunded, then that
forwarder's long-term key was compromised.
\end{lemma}
\begin{spthy}
lemma Intermediary_Never_Loses_Under_Liveness:
  "All P Prev Next Up ptrIn ptrOut y #i #j #k.
      HTLCForwarded(P, Prev, Next, ptrIn, ptrOut, y) @ #i
    & Redeemed(ptrOut, y) @ #j
    & Refunded(Up, ptrIn, y) @ #k
   ==> Ex #l. LtkCompromised(P) @ #l"
\end{spthy}
Lemma~\ref{lem:intermediary} does not say when the key was compromised: a compromise after the
refund also satisfies it. The strong form closes this gap.

\begin{lemma}[Intermediary Safety, Strong Form]\label{lem:intermediary-strong}
If a forwarder's outgoing HTLC is redeemed while its incoming HTLC is refunded, then that
forwarder's key was compromised before its incoming HTLC timed out.
\end{lemma}
\begin{spthy}
lemma Intermediary_Never_Loses_Strong:
  "All P Prev Next Up ptrIn ptrOut y #i #j #k.
      HTLCForwarded(P, Prev, Next, ptrIn, ptrOut, y) @ #i
    & Redeemed(ptrOut, y) @ #j
    & Refunded(Up, ptrIn, y) @ #k
   ==> Ex #l #t. LtkCompromised(P) @ #l & TimedOut(ptrIn) @ #t & #l < #t"
\end{spthy}
Lemma~\ref{lem:intermediary-strong} matches the honesty condition of T2 and T3, which bind a
forwarder only if it is not compromised before its incoming HTLC times out. Both lemmas
quantify over every forward, so they hold for every hop of a route of any length, in every
trace that satisfies T1, T2 and T3.

Routing causality says that nothing is forwarded out of thin air.

\begin{lemma}[Routing Causality]\label{lem:causality}
Every forward on hash $y$ is preceded by the HTLC it forwards: either the sender's offer on the
forwarder's incoming channel, or the predecessor's forward that locked that channel's HTLC.
\end{lemma}
\begin{spthy}
lemma Forward_Requires_Incoming [use_induction]:
  "All Me Prev Next ptrIn ptrOut y #i.
      HTLCForwarded(Me, Prev, Next, ptrIn, ptrOut, y) @ #i
   ==> (Ex S #j. HTLCOffered(S, Me, ptrIn, y) @ #j & #j < #i)
     | (Ex P0 pin0 #j. HTLCForwarded(Prev, P0, Me, pin0, ptrIn, y) @ #j & #j < #i)"
\end{spthy}
Because \code{lightning.spthy} has a single generic \code{Forward_HTLC} rule, this one lemma
covers every hop of a route of any length, where a model with a fixed path needs one lemma per
position (and an honest variant of each). It needs no honesty premise, because a forward
consumes the linear \code{Route} fact, which only \code{Sender_Offer_HTLC} and
\code{Forward_HTLC} produce. On the release side, an HTLC is redeemed only by the receiver's
fulfilment or by a forwarder settling that hop (\code{Redeemed_Source}), and every preimage
that a forwarder or the sender takes from a channel in order to settle comes from an earlier
redemption of the HTLC on that channel or from an earlier wormhole skip
(\code{Fulfilled_Source}, Sec.~\ref{sec:wormhole}).

Finally, a forwarder is paid for its service only if the service succeeded.

\begin{lemma}[Fees Only on Success]\label{lem:fee}
A forwarder whose key is never compromised earns the fee of a hop only after it has forwarded
on that hop and its outgoing HTLC has been redeemed.
\end{lemma}
\begin{spthy}
lemma Fee_Only_On_Successful_Forward [heuristic={route_info_first}]:
  "All Me ptrIn ptrOut y %fee #i.
      FeeEarned(Me, ptrIn, ptrOut, y, %fee) @ #i
    & not (Ex #c. LtkCompromised(Me) @ #c)
   ==> (Ex #j. Redeemed(ptrOut, y) @ #j & #j < #i)
     & (Ex Prev Next #f. HTLCForwarded(Me, Prev, Next, ptrIn, ptrOut, y) @ #f & #f < #i)"
\end{spthy}
Per hop, the incoming amount is the outgoing amount plus the fee by construction (the \code{Eq}
action of \code{Forward_HTLC}, Sec.~\ref{sec:core-amounts}). Amounts are not chained along the
route, however: the offer, the invoice and the \code{Route} fact carry no amount, so we make no
end-to-end value claim.

% The label sits before \paragraph so that it prints as "V-C" (IEEEtran would print "V-C0a").
\phantomsection\label{par:nonvacuity}%
\paragraph{Non-vacuity}
Safety lemmas are useless if the model cannot run at all, so we also checked
\code{exists-trace} witnesses. \code{Multihop_Routing_Reachable} shows an HTLC routed from the
sender over two forwarders to the receiver, who fulfils it.
\code{Sanity_Honest_Payment_Completes} and \code{Sanity_Atomic_Payment_Completes} show payments
over one and over two forwarders that settle end to end, with no key compromise and no refund:
every forwarder claims its incoming HTLC and the sender settles. The second witness satisfies
the premises of Lemma~\ref{lem:atomicity}. \code{Distinct_Parties_Configuration} shows three
channels among four different parties, with no key compromise and no dispute.
\code{Sanity_Compromised_Intermediary_Can_Lose} shows a compromised forwarder whose outgoing
HTLC is redeemed while its incoming HTLC is refunded, so the premises of
Lemmas~\ref{lem:intermediary} and~\ref{lem:intermediary-strong} are satisfiable.

\subsection{Timing and Liveness}\label{sec:timing}

Our model has three liveness restrictions (Sec.~\ref{sec:restrictions}): T1 is
\code{RedeemBeforeTimeout}, T2 is \code{OutgoingTimesOutFirst} and T3 is
\code{IntermediaryMustClaim}. The order of a hop's two timeouts is not assumed separately:
\code{ForwardTimeGap} is a lemma, derived from T2 and verified.

Concrete block heights are treated in two auxiliary models of \code{lightning_aux.spthy}. With
\code{-D=CLTV}, a positive CLTV delta (restriction \code{DeltaPositive}) leaves every forwarder
a non-empty claim window: a preimage revealed before the outgoing deadline is revealed before
the incoming one (\code{Claim_Window_Nonempty}), also over two staggered hops
(\code{Staggered_Path_Safe}). With \code{-D=CLOCK}, channels carry concrete deadlines on a block
clock, and nine lemmas (six all-traces, three exists-trace) cover the channel and HTLC
lifecycle, e.g., \code{No_HTLC_After_Close}; we return to two of them below.

The four-rule model \code{lightning_aux.spthy} \code{-D=TIMEOUT} has no ordering restriction at
all. There \tamarin{} finds the early-timeout race (\code{Early_Timeout_Race}, exists-trace,
6 steps): a forwarder's incoming HTLC times out and is refunded, and its outgoing HTLC is
redeemed after that timeout. In \code{lightning.spthy}, the race is excluded for every
forwarder that is not compromised before its incoming timeout:
\begin{spthy}
lemma Timeout_Race_Blocked:
  "not (Ex P Prev Next Up ptrIn ptrOut y #l #i #j #k.
      HTLCForwarded(P, Prev, Next, ptrIn, ptrOut, y) @ #l
    & TimedOut(ptrIn)     @ #k
    & Refunded(Up, ptrIn, y)  @ #j
    & Redeemed(ptrOut, y) @ #i
    & #k < #i
    & not (Ex #c. LtkCompromised(P) @ #c & #c < #k))"
\end{spthy}
This is a two-sided check: the race is reachable without ordering restrictions and blocked in
\code{lightning.spthy}. \code{Timeout_Race_Blocked} stays verified when any single one of T1,
T2 and T3 is removed (Table~\ref{tab:necessity}).

\begin{table}[t]
  \caption{Necessity of the liveness restrictions (\texttt{make necessity}:
    \texttt{lightning.spthy} with \texttt{-D=NECESSITY} and one restriction removed per column,
    600\,s per check). verified~($n$): proved in $n$ steps; no result: neither a proof nor a
    counterexample within 600\,s, from which we draw no conclusion. The counterexample
    \texttt{Necessity\_IMC\_Honest\_Intermediary\_Loses} is checked only with T3 removed.}
  \label{tab:necessity}
  \centering
  \footnotesize
  \setlength{\tabcolsep}{3pt}
  \renewcommand{\arraystretch}{1.12}
  \begin{tabular}{@{}L{3.45cm}ccc@{}}
    \toprule
    Lemma & T1 removed & T2 removed & T3 removed\\
          & {\scriptsize\code{-D=NO_RBT}} & {\scriptsize\code{-D=NO_OTF}}
          & {\scriptsize\code{-D=NO_IMC}}\\
    \midrule
    \code{ForwardTimeGap} & verified (6) & no result & verified (6)\\
    \code{Timeout_Race_Blocked} & verified (4) & verified (4) & verified (4)\\
    \code{Intermediary_Never_Loses_Under_Liveness} (Lemma~\ref{lem:intermediary})
      & verified (4) & verified (4) & no result\\
    \code{Intermediary_Never_Loses_Strong} (Lemma~\ref{lem:intermediary-strong})
      & verified (4) & verified (4) & no result\\
    \code{Payment_Atomicity_Under_Liveness} (Lemma~\ref{lem:atomicity})
      & verified (10) & no result & no result\\
    \midrule
    \code{Necessity_IMC_Honest_Intermediary_Loses} (exists-trace)
      & -- & -- & verified (28)\\
    \bottomrule
  \end{tabular}
\end{table}

We removed each liveness restriction in turn (\texttt{make necessity}) and re-ran the five
lemmas that may depend on them (Table~\ref{tab:necessity}). Without T3, \tamarin{} finds a counterexample
(exists-trace, 28 steps):
\begin{spthy}
lemma Necessity_IMC_Honest_Intermediary_Loses:
  exists-trace
  "Ex S F1 R p0 p1 x #o #a #r #u.
      HTLCOffered(S, F1, p0, h(x)) @ #o
    & HTLCForwarded(F1, S, R, p0, p1, h(x)) @ #a
    & HTLCFulfilled(R, F1, p1, h(x)) @ #r
    & Refunded(S, p0, h(x)) @ #u
    & not (Ex P #k. LtkCompromised(P) @ #k)
    & not (Ex ptr y1 y2 #m #n. HTLCAdded(ptr, y1) @ #m & HTLCAdded(ptr, y2) @ #n & not (#m = #n))
    & not (Ex Me P q z #e. IntermediaryClaimed(Me, P, q, z) @ #e)
    & not (Ex A B C q z #e. SenderSettled(A, B, C, q, z) @ #e)"
\end{spthy}
No key is compromised, \code{F1} forwards, \code{R} redeems the outgoing HTLC of \code{F1},
\code{F1} never claims its incoming HTLC, and the HTLC of \code{S} to \code{F1} is refunded: the
honest intermediary \code{F1} loses, which violates Lemmas~\ref{lem:intermediary}
and~\ref{lem:intermediary-strong}. (The direct proof attempts of these two lemmas without T3
gave no result within 600\,s; the counterexample is this separate exists-trace lemma.) T2 only
orders the two timeouts; it does not force the intermediary to claim its incoming HTLC. That
claim is exactly T3, and no ordering rule can force an action. So T3 is load-bearing for
intermediary safety.

The two ordering restrictions are not needed in this way. T1 is not needed by any of the five
lemmas: all of them stay verified without it. T2 is not needed for intermediary safety:
Lemmas~\ref{lem:intermediary} and~\ref{lem:intermediary-strong} stay verified without it. We
draw no conclusion from the cells of Table~\ref{tab:necessity} without a result; in particular,
we do not claim that \code{ForwardTimeGap} or Lemma~\ref{lem:atomicity} needs T2, or that
Lemma~\ref{lem:atomicity} needs T3.

% The label sits before \paragraph so that it prints as "V-D" (IEEEtran would print "V-D0a").
\phantomsection\label{par:liveness-scope}%
\paragraph{Liveness as Restrictions}
Two of our three liveness assumptions, T2 and T3, are \tamarin{} restrictions of the form
\emph{if A happens, then B happens} (T1 only orders a redemption before a timeout); in T3, for
instance, B is the redemption of the forwarder's incoming HTLC. A restriction removes every
trace that violates it, including traces in which B is impossible rather than merely late. A
lemma proved under T3 therefore says nothing about runs in which an honest forwarder
\emph{cannot} claim in time, which is where real timing attacks live.
Likewise, T2 removes every trace in which an honest forwarder's incoming HTLC times out while
its outgoing HTLC has not timed out earlier. We make the scope of this assumption explicit.

\emph{Which results depend on it.} Removing T3 breaks intermediary safety (the counterexample
above); Table~\ref{tab:necessity} shows which of the five liveness lemmas are still proved
without each restriction. The 24 lemmas of \texttt{make general},
namely the bypass theorem for routes of any length (Sec.~\ref{sec:bypass-any}), its helpers and
16 lemmas of \texttt{make verify}, among them \code{Skipped_Requires_Compromise} and
\code{Refund_Requires_Timeout}, are proved with all three liveness restrictions removed. The
wormhole trace \code{Wormhole_Fee_Theft_Reachable} is verified with every restriction in place,
so the restrictions do not hide it. All other all-traces lemmas of \texttt{make verify},
including the bypass theorems for positions 1 to 6, are proved with all three restrictions in
place; apart from the five lemmas of Table~\ref{tab:necessity}, we have not checked which of
them still hold without them.

\emph{A claim window exists, in a separate model.} The Clock model
(\code{lightning_aux.spthy} \code{-D=CLOCK}) has no liveness restriction. It proves
\code{Claim_Window_Exists}: if a forwarder's outgoing HTLC is redeemed and its incoming HTLC is
refunded, the redemption happens at a strictly lower block height than the refund, so the
forwarder had a non-empty window in which to claim. It also proves
\code{Loss_Implies_Skipped_Claim}: a forwarder that loses in this way never claimed its incoming
HTLC. The Clock model is not composed with \code{lightning.spthy}, so it supports T3 but does
not discharge it.

\emph{Alternatives.} A more faithful treatment would replace the restrictions by the progress
semantics of Backes et al.~\cite{backes2017}. Roughly, honest processes that can continue
without input from the adversary eventually do, so traces in which the forwarder cannot act
are kept instead of pruned. \tamarin{} 1.10 implements this for \sapic{} processes (option
\code{translation-progress}), so the routing layer would have to be written as \sapic{}
processes. A cheaper first step is to turn T1, T2 and T3 into premises of only the lemmas that
need them, so that they no longer prune traces for every other lemma. We leave both to future
work (Sec.~\ref{sec:limitations}).

\subsection{Griefing and the Wormhole}\label{sec:wormhole}

\code{Refund_Possible} shows a griefing trace: the sender offers an HTLC against an invoice,
nobody forwards it, and it is refunded, so the payment never settles; no key is compromised.
Since a refund requires a timeout (\code{Refund_Requires_Timeout}), the capital stays locked
until then. This witness refunds the sender, not an intermediary. We report it as reachable
griefing (denial of service), not as a new attack; it belongs to group 6 of
Table~\ref{tab:groups}, the circled 6 in Fig.~\ref{fig:griefing}.

\begin{figure}[t]
  \centering
  \definecolor{wgreen}{RGB}{27,94,32}
  \definecolor{wred}{RGB}{183,28,28}
  \definecolor{worange}{RGB}{200,90,0}
  \begin{tikzpicture}[y=0.88cm,
      font=\sffamily\footnotesize,
      >={Stealth[length=4pt,width=3.5pt]},
      party/.style={rounded corners=2pt, minimum width=1.38cm, minimum height=0.5cm,
                    inner sep=1pt, text=white, font=\sffamily\footnotesize\bfseries},
      honest/.style={party, fill=wgreen},
      attacker/.style={party, fill=wred},
      life/.style={draw=black!30, dashed, thin},
      msg/.style={->, draw=wgreen, semithick},
      lbl/.style={font=\sffamily\scriptsize, inner sep=1.5pt, fill=white},
    ]
    % parties
    \node[honest]   (A)  at (0.0, 0) {Alice};
    \node[attacker] (M1) at (1.8, 0) {Att1};
    \node[honest]   (I)  at (3.6, 0) {Indiana};
    \node[attacker] (M2) at (5.4, 0) {Att2};
    \node[honest]   (B)  at (7.2, 0) {Bob};
    \foreach \p/\x in {A/0.0, M1/1.8, I/3.6, M2/5.4, B/7.2}
      \draw[life] (\p.south) -- (\x, -4.2);
    % invoice
    \draw[->, draw=wgreen, semithick, dashed] (7.2, -0.7) -- (0.0, -0.7)
      node[lbl, above, pos=0.5, text=wgreen] {invoice $y = h(x)$};
    % lock: HTLC chain
    \draw[msg] (0.0, -1.12) -- (1.8, -1.12) node[lbl, above, midway] {HTLC($y$)};
    \draw[msg] (1.8, -1.45) -- (3.6, -1.45)  node[lbl, above, midway] {HTLC($y$)};
    \draw[msg] (3.6, -1.85) -- (5.4, -1.85) node[lbl, above, midway] {HTLC($y$)};
    \draw[msg] (5.4, -2.25) -- (7.2, -2.25)  node[lbl, above, midway] {HTLC($y$)};
    % release: Bob redeems Att2's HTLC
    \draw[msg] (7.2, -2.75) -- (5.4, -2.75) node[lbl, above, midway] {redeem($x$)};
    % off-path relay around Indiana
    \draw[->, draw=wred, thick, dashed] (5.4, -2.95) .. controls (4.6, -3.75) and (2.6, -3.75)
      .. (1.8, -2.95);
    \node[lbl, text=wred, align=center] at (3.6, -3.85)
      {$x$ relayed off-path (\textsf{AttackerSkip})};
    \node[circle, fill=wred, text=white, inner sep=0.6pt, minimum size=9pt,
          font=\sffamily\scriptsize\bfseries] at (1.25, -3.85) {7};
    % Att1 claims Alice's HTLC
    \draw[msg] (1.8, -3.3) -- (0.0, -3.3) node[lbl, above, midway] {redeem($x$)};
    % outcomes
    \node[lbl, text=wgreen, font=\sffamily\scriptsize\bfseries] at (0.0, -4.45) {settles};
    \node[lbl, text=wgreen, font=\sffamily\scriptsize\bfseries] at (7.2, -4.45) {is paid};
    \node[lbl, draw=worange, rounded corners=2pt, text=worange, align=center,
          font=\sffamily\scriptsize\bfseries, inner sep=2pt] at (3.6, -4.5)
      {refunded on timeout,\\ earns no fee};
  \end{tikzpicture}
  \caption{Wormhole (\texttt{Wormhole\_Fee\_Theft\_Reachable}). Att1 and Att2 are compromised
    (\texttt{AttackerSkip} requires this) and not adjacent; Alice, Indiana and Bob are honest.
    Att2 relays $x$ to Att1 around Indiana, so Alice settles and Bob is paid, while Indiana's
    outgoing HTLC is refunded and Indiana earns no fee. The lemma fixes no order of its events;
    we draw them in causal order.}
  \label{fig:wormhole}
\end{figure}

The wormhole attack~\cite{malavolta2019} is an instance of the agent-skipping attack formalised
symbolically by~\cite{smith2022}. It is a known attack, so we do not claim it as new. First,
after the linear single-spend fix (Sec.~\ref{par:single-spend}), it is a soundness check: the
model must still admit this behaviour, or the fix would be too strong. Second, the model shows
who can mount the attack. The wormhole is the attacker rule \code{AttackerSkip}
(Sec.~\ref{sec:core-skip}): two compromised parties on the same payment pass the preimage $x$
from the exit party's outgoing HTLC to the entry party's outgoing HTLC, around the hops between
them.

\code{Wormhole_Fee_Theft_Reachable} (exists-trace, verified by replaying a saved 116-step proof)
is the five-party trace of Fig.~\ref{fig:wormhole}, marked by the circled 7. Bob creates an
invoice for Alice, and the HTLC on $y = h(x)$ is locked hop by hop from Alice through Att1,
Indiana and Att2 to Bob. Bob redeems the HTLC of Att2 with $x$ and is paid. \code{AttackerSkip}
then relays $x$ from Att2 to Att1, off the path, and Att1 claims the HTLC of Alice, so Alice
settles. Indiana's outgoing HTLC to Att2 is refunded on timeout, and Indiana earns no fee. The
lemma requires that Alice, Bob and Indiana are never compromised; in the trace, Att1 and Att2
are compromised, as \code{AttackerSkip} requires. An honest hop is thus bypassed although the
payment settles (full statement in Appendix~\ref{app:catalogue}). Conversely, a skip requires a
key compromise:

\begin{lemma}[Skips Require Compromise]\label{lem:skip}
Every wormhole skip happens after some party's long-term key was compromised.
\end{lemma}
\begin{spthy}
lemma Skipped_Requires_Compromise [reuse]:
  "All p x #i. Skipped(p, x) @ #i ==> Ex P #c. LtkCompromised(P) @ #c & #c < #i"
\end{spthy}
Lemma~\ref{lem:skip} is also re-proved without the liveness restrictions (\texttt{make general}).
It does not exclude the wormhole; it says that the wormhole needs a compromise.
Sec.~\ref{sec:bypass-fixed} proves, for the forwarder at each position 1 to 6 of a route
(Indiana is at position 2), that if the sender settles while that forwarder's outgoing HTLC is
refunded, some party's key was compromised (group 9 of Table~\ref{tab:groups}; under T1 to T3,
with no assumption on the number of offers). Sec.~\ref{sec:bypass-any} extends this to routes
of any length, for payments whose hash is offered once, without the liveness restrictions.
Neither result says which party was compromised, or when. Sec.~\ref{sec:ptlc} shows, on a fixed
route of \code{lightning_aux.spthy}, that per-hop locks rule out the attack around an honest hop
when the sender is honest (\code{No_Wormhole_Around_Honest_Hop}), whereas a shared hash lock
admits it (\code{Wormhole_Around_Honest_Hop}).
